\SourcePage{81}{84}
\section{The Variational Principle}

The Schrödinger equation in its general form, $\hat H\psi=E\psi$, can be
obtained from the variational principle
\begin{equation}
  \delta\int\psi^*(\hat H-E)\psi\,dq=0.
  \tag{20.1}
\end{equation}
Since $\psi$ is complex, variations with respect to $\psi$ and $\psi^*$ may
be made independently. Variation with respect to $\psi^*$ gives
\[
  \int\delta\psi^*(\hat H-E)\psi\,dq=0,
\]
and the arbitrariness of $\delta\psi^*$ then yields the required equation
$\hat H\psi=E\psi$. Variation with respect to $\psi$ gives nothing new.
Indeed, varying $\psi$ and using the Hermitian character%
\SourcePage{82}{85}
of $\hat H$, we have
\[
  \int\psi^*(\hat H-E)\delta\psi\,dq
  =\int\delta\psi(\hat H^*-E)\psi^*\,dq=0,
\]
which yields the complex-conjugate equation
$\hat H^*\psi^*=E\psi^*$.

The variational principle (20.1) requires an unconstrained extremum of the
integral. It may be put in another form by treating $E$ as a Lagrange
multiplier in the problem of finding a constrained extremum of
\begin{equation}
  \delta\int\psi^*\hat H\psi\,dq=0
  \tag{20.2}
\end{equation}
subject to the supplementary condition
\begin{equation}
  \int\psi\psi^*\,dq=1.
  \tag{20.3}
\end{equation}

The minimum value of the integral (20.2), subject to (20.3), is the first
energy eigenvalue, that is, the energy $E_0$ of the normal state. The function
$\psi$ which realizes this minimum is correspondingly the normal-state wave
function $\psi_0$\footnote{For the remainder of this section we take the wave
functions $\psi$ to be real, as they can always be chosen when no magnetic
field is present.}. The wave functions $\psi_n$ ($n>0$) of the succeeding
stationary states correspond only to extrema, not to true minima of the
integral.

To obtain from the minimum condition for (20.2) the wave function $\psi_1$
and energy $E_1$ of the state immediately above the normal state, the trial
functions $\psi$ must be restricted not only by the normalization condition
(20.3), but also by orthogonality to the normal-state wave function $\psi_0$:
$\int\psi\psi_0\,dq=0$. In general, suppose that the wave functions
$\psi_0,\psi_1,\ldots,\psi_{n-1}$ of the first $n$ states are known (the
states being ordered by increasing energy). The wave function of the next
state minimizes (20.2) under the supplementary conditions
\begin{equation}
  \int\psi^2\,dq=1,
  \qquad
  \int\psi\psi_m\,dq=0,
  \qquad
  m=0,1,2,\ldots,n-1.
  \tag{20.4}
\end{equation}
\SourcePage{83}{86}
We give here several general theorems which may be proved from the variational
principle\footnote{Proofs of the theorems about zeros of eigenfunctions (see
also the next section) may be found in: \emph{M.~A.~Lavrent'ev and
L.~A.~Lyusternik}, Course of the Calculus of Variations, Moscow:
Gostekhizdat, 1950, Ch.~IX; \emph{R.~Courant and D.~Hilbert}, Methods of
Mathematical Physics, Moscow: Gostekhizdat, 1951, Vol.~I, Ch.~VI.}.

The normal-state wave function $\psi_0$ does not vanish (or, as one says, has
no nodes) at any finite values of the coordinates\footnote{This theorem (and
the consequences drawn from it below) is not valid in general for wave
functions of systems containing several identical particles; see the end of
\secref{63}.}. In other words, it has the same sign throughout space. It follows
that the wave functions $\psi_n$ ($n>0$) of the other stationary states,
being orthogonal to $\psi_0$, necessarily have nodal points (if $\psi_n$ also
had a constant sign, the integral $\int\psi_0\psi_n\,dq$ could not vanish).

The absence of nodes in $\psi_0$ also implies that the normal energy level
cannot be degenerate. Suppose otherwise, and let $\psi_0$ and $\psi'_0$ be
two different eigenfunctions belonging to the level $E_0$. Every linear
combination $c\psi_0+c'\psi'_0$ would also be an eigenfunction; but by a
suitable choice of $c$ and $c'$, this function can always be made to vanish
at any prescribed point of space, producing an eigenfunction with a node.

If motion is confined to a bounded region of space, then $\psi=0$ on its
boundary (see \secref{18}). To determine the energy levels one must use the
variational principle to find the minimum of (20.2) with this boundary
condition. In this case the theorem that the normal-state wave function has
no nodes states that $\psi_0$ vanishes nowhere inside the region.

Notice that increasing the size of the region of motion lowers every energy
level $E_n$. This follows directly because enlarging the region enlarges the
class of trial functions available for minimizing the integral, so its
minimum value can only decrease. For discrete-spectrum states of a system of
particles, the expression
\[
  \int\psi\hat H\psi\,dq
  =\int\left[-\sum_a\frac{\hbar^2}{2m_a}\psi\Delta_a\psi
  +U\psi^2\right]dq
\]
can be%
\SourcePage{84}{87}
transformed into a form more convenient for actual variation. In the first
term of the integrand, write
\[
  \psi\Delta_a\psi
  =\DivAt{a}{(\psi\GradAt{a}{\psi})}
  -{\left(\GradAt{a}{\psi}\right)}^2.
\]
The integral of $\DivAt{a}{(\psi\GradAt{a}{\psi})}$ over $dV_a$ becomes an integral
over a closed surface at infinity. Since discrete-spectrum wave functions
vanish sufficiently rapidly at infinity, this integral disappears. Thus
\begin{equation}
  \int\psi\hat H\psi\,dq
  =\int\left[\sum_a\frac{\hbar^2}{2m_a}
  {\left(\GradAt{a}{\psi}\right)}^2
  +U\psi^2\right]dq.
  \tag{20.5}
\end{equation}
